% BHU PYQ -- Professional Exam Template
\documentclass[11pt,a4paper]{article}

\usepackage[a4paper,top=2.5cm,bottom=2.5cm,left=2.8cm,right=2.8cm,headheight=14pt]{geometry}
\usepackage{amsmath,amssymb}
\usepackage[T1]{fontenc}
\usepackage[utf8]{inputenc}
\usepackage{lmodern}
\usepackage{microtype}
\usepackage{fancyhdr}
\usepackage{enumitem}
\usepackage{listings}
\usepackage{hyperref}
\usepackage{lastpage}
\usepackage{parskip}

\hypersetup{colorlinks=true,urlcolor=black,linkcolor=black,
  pdftitle={BVCA-205: Skill Training Python -- BHU PYQ}}

\pagestyle{fancy}
\fancyhf{}
\renewcommand{\headrulewidth}{0.4pt}
\renewcommand{\footrulewidth}{0.4pt}
\fancyhead[L]{\small   \textit{Banaras Hindu University}}
\fancyhead[R]{\small   \textit{BVCA-205: Skill Training Python}}
\fancyfoot[C]{\small Page  \thepage\ of \pageref{LastPage}}

\lstset{
  basicstyle=\small tfamily,
  language=Python,
  frame=single,
  breaklines=true,
  columns=flexible,
  keepspaces=true
}


\newlist{parts}{enumerate}{2}
\setlist[parts,1]{label=(\alph*),leftmargin=*,itemsep=4pt,topsep=3pt}
\setlist[parts,2]{label=(\roman*),leftmargin=*,itemsep=2pt,topsep=2pt}

\newcommand{\pts}[1]{\hfill   \textit{[#1]}}


\begin{document}


\begin{center}
  {\large\bfseries BANARAS HINDU UNIVERSITY}\\[2pt]
  {\normalsize B.Voc. Semester II Examination, 2021-22}\\[6pt]
  \rule{0.8\linewidth}{0.5pt}\\[6pt]
  {\large\bfseries Paper: BVCA-205 --- Skill Training in Python}\\[4pt]
  {\normalsize Time: Three Hours \hfill Maximum Marks: 70}
\end{center}

\rule{\linewidth}{0.4pt}

\smallskip



\noindent   \textit{Instructions: Attempt all questions. Part A: Multiple Choice and Short Questions. Part B: Internal Choice (any four). Part C: Open choice (any three out of four).}

\bigskip

\bigskip


\noindent {\large\bfseries PART --- A}
\rule{\linewidth}{0.4pt}\smallskip




\noindent   \textit{Note: Multiple Choice Type Questions and Short Questions.} \hfill   \textbf{[10 $ \times$ 2 = 20]}

\medskip


\noindent   \textbf{Question 1.} Which of the following is a Python tuple?

\begin{parts}
  \item   \texttt{\{1, 2, 3\}}
  \item   \texttt{4}
  \item   \texttt{[1, 2, 3]}
  \item   \texttt{(2, 2, 3)}
\end{parts}

\medskip


\noindent   \textbf{Question 2.} Which one of the following is not a keyword in Python?

\begin{parts}
  \item   \texttt{pass}
  \item   \texttt{eval}
  \item   \texttt{assert}
  \item   \texttt{nonlocal}
\end{parts}

\medskip


\noindent   \textbf{Question 3.} In object-oriented programming, what is the relation of an object to a class?

\begin{parts}
  \item Class is an instance of an object.
  \item Object is a child of a class.
  \item Object is a method of a class.
  \item Object is an instance of a class.
\end{parts}

\medskip


\noindent   \textbf{Question 4.} What will be the output of the following Python code snippet?

\begin{lstlisting}
print(type(5/2))
print(type(5//2))
\end{lstlisting}

\begin{parts}
  \item float and int
  \item float and float
  \item int and int
  \item int and float
\end{parts}

\medskip


\noindent   \textbf{Question 5.} What will be the output of the following Python code snippet?

\begin{lstlisting}
print(2**4 + (5+5)**(1+1))
\end{lstlisting}

\begin{parts}
  \item Error
  \item 28
  \item 118
  \item 116
\end{parts}

\medskip


\noindent   \textbf{Question 6.} What will be the output of the following Python code?

\begin{lstlisting}
i = 1
while True:
    if i % 3 == 0:
        break
    print(i)
    i += 1
\end{lstlisting}

\begin{parts}
  \item 1 2 3
  \item Error
  \item 1 2
  \item None of the mentioned
\end{parts}

\medskip


\noindent   \textbf{Question 7.} Which one of the following has the same precedence level?

\begin{parts}
  \item Addition and Subtraction
  \item Multiplication, Division and Addition
  \item Multiplication, Division, Addition and Subtraction
  \item Addition and Multiplication
\end{parts}

\medskip


\noindent   \textbf{Question 8.} What will be the datatype of   \texttt{test} in the below code snippet?

\begin{lstlisting}
test = 10
print(type(test))
test = "testbook"
print(type(test))
\end{lstlisting}

\begin{parts}
  \item str and int
  \item int and int
  \item str and str
  \item int and str
\end{parts}

\medskip


\noindent   \textbf{Question 9.} What is the meaning of   \texttt{>>>} in Python?

\begin{parts}
  \item Interpreter is ready to take instruction.
  \item Compiler is ready to take instruction.
  \item Three right shift.
  \item Three left shift.
\end{parts}

\medskip


\noindent   \textbf{Question 10.} What is the output of   \texttt{print(math.pow(3, 2))}?

\begin{parts}
  \item 9.0
  \item None
  \item 9
  \item None of the mentioned
\end{parts}

\bigskip


\noindent {\large\bfseries PART --- B}
\rule{\linewidth}{0.4pt}\smallskip




\noindent   \textit{Note: Internal Choice --- Either/Or type questions. Each question carries 5 marks. Attempt any four.} \hfill   \textbf{[5 $ \times$ 4 = 20]}

\medskip


\noindent   \textbf{Question 11.} What are the common built-in data types in Python?
\centerline{   \textbf{OR}}
How do you read and write into Python files? Explain with a program.

\medskip


\noindent   \textbf{Question 12.} Explain the concept of overriding methods with a program.
\centerline{   \textbf{OR}}
Write a Python program to create a lambda function that returns the square value of a given number.

\medskip


\noindent   \textbf{Question 13.} Write a Python program to create a static method that calculates the square root value of a given number.
\centerline{   \textbf{OR}}
Explain the   \texttt{values()} and   \texttt{items()} methods used in dictionaries with examples.

\medskip


\noindent   \textbf{Question 14.} What is the difference between Python Arrays, lists, and tuples with proper examples?
\centerline{   \textbf{OR}}
Write a Python program to define a Student class and create an object. Also, display the student's details.

\medskip


\noindent   \textbf{Question 15.} What are the key features of Python over other programming languages?
\centerline{   \textbf{OR}}
Write a function that accepts two values and finds their sum.

\bigskip


\noindent {\large\bfseries PART --- C}
\rule{\linewidth}{0.4pt}\smallskip




\noindent   \textit{Note: Open choice --- Attempt any 3 out of 4 questions. Each question carries 10 marks.} \hfill   \textbf{[3 $ \times$ 10 = 30]}

\medskip


\noindent   \textbf{Question 16.} List the operators that Python supports. Explain the relational and logical operators along with their precedence while evaluating an expression.

\medskip


\noindent   \textbf{Question 17.} Write short notes on:

\begin{parts}
  \item Debugging
  \item Exception Handling
\end{parts}

\medskip


\noindent   \textbf{Question 18.} Write a Python program to retrieve the elements from a 2D array and display them using two different loops.

\medskip


\noindent   \textbf{Question 19.} Discuss the technical strengths of Python. What is the difference between list and tuples in Python?

\vfill
\rule{\linewidth}{0.4pt}

\begin{center}\small
\href{https://bhu-exam.vercel.app/}{Click here to visit our website}\\[4pt]\href{https://me-aryan.vercel.app/}{Click here to visit our contributor}\\[4pt]\href{https://quashan-me.vercel.app/}{Click here to visit our contributor}
\end{center}

\end{document}
